\subsection{FRM EXAM 2009---QUESTION 5-14: which statement about ES is incorrect?}

\textbf{Enunciado}

\emph{Which of the following statements about expected shortfall (ES) is incorrect?}
\begin{enumerate}[label=\alph*.]
\item ES provides a consistent risk measure across different positions and takes account of correlations.
\item ES tells what to expect in bad states: it gives an idea of how bad the payoff can be expected to be if the outcome is bad.
\item ES-based rule is consistent with expected utility maximization if risks are ranked by a second-order stochastic dominance rule.
\item \hl{Like VAR, ES does not always satisfy subadditivity (i.e., the risk of a portfolio must be less than or equal to the sum of the risks of its individual positions).}
\end{enumerate}

\subsubsection*{Solución}


\textbf{Propiedad de coherencia. }El Expected Shortfall es una medida de riesgo \textbf{coherente}, por lo que satisface subaditividad:
\[
\ES_\alpha(L_A+L_B)\le \ES_\alpha(L_A)+\ES_\alpha(L_B).
\]
En contraste, el VaR puede violar subaditividad.


El inciso (d) afirma que ES ``no siempre satisface subaditividad'' como VaR; eso es falso.
Así, (d) es el inciso incorrecto.












